\documentclass[11pt,a4paper,openright,oneside]{report}
\usepackage[utf8]{inputenc}

\title{\textbf{Aula3}}
\author{Rafael Pinto}
\date{20/10/2020}

\begin{document}

Olá 
Finalmente, outro pacote que é útil para português é ofontenc, que indica o tipo decodificação de carateres do resultado da compilação do LATEX. Aconselha-se o uso dacodificaçãoT1, porque a mesma permite lidar convenientemente com carateres acentuados.

\part{ola}
\chapter{tugj}
\tableofcontents
\section{bola}
\section{oooo}
\subsection{e}
\subsubsection{titi}
\renewcommand{\partname}{Parte}

\begin{table}
\caption{123}
\centerline{456}
\label{feito}
\end{table}

\begin{figure}[h]
\centerline{\fbox{ronaldo}}
\caption{cr7}
\end{figure}

\includegraphics[height=24pt]{/home/rafael/labi/índice.png}
\end{document}
